%% notes-tex/025-additive-baselines/sections/6-disjoint-transformer.tex

\section{The transformer on held-out screens\secstatus{todo}}
\label{sec:disjoint}

\subsection{The run registry\secstatus{todo}}
\label{sec:registry}

Every transformer number in this document traces to one of eight training runs.
Table~\ref{tab:wandbruns} names each once, with the configuration it ran and what
that configuration turns on, and links its W\&B page, which holds the per-epoch
training loss, validation Pearson and the selected epoch. The table is generated
from the W\&B API and no run id is typed into the text.

%% SOURCE: experiments/025-solid-growth/scripts/wandb_run_index_025.py ->
%% tables/t5-wandb-runs.tex and results/wandb_run_index_025.json; every row is
%% verified against the W&B API when the script runs.
\begin{table}[htbp]
\centering
\scriptsize
\caption{The run registry. The 010 checkpoints are arm R's reference and were scored
on the 010 test split by their own evaluation runs. GH 1598 and GH 1640 were stopped
by the 12 hour wall clock and have no test evaluation; GH 1640's epoch 7 checkpoint
was scored afterward on CPU. The IGB runs finished 30 epochs with no test
evaluation. The validation maximum is an upward-biased order statistic over the
logged epochs. Configs of the IGB runs are on branch
\texttt{feat/025-fitness-joint-head}.}
\label{tab:wandbruns}
%% GENERATED by experiments/025-solid-growth/scripts/wandb_run_index_025.py
%% SOURCE: results/wandb_run_index_025.json, each row verified against the W&B API. Do not edit by hand.
\begin{tabular}{@{}llll>{\raggedright\arraybackslash}p{19mm}>{\raggedright\arraybackslash}p{24mm}rr>{\raggedright\arraybackslash}p{15mm}@{}}
\toprule
Run & W\&B & Arm & Config & Gene input & Schedule, readout & Epochs & Val max (ep.) & Test \\
\midrule
CGT M01 & \href{https://wandb.ai/zhao-group/torchcell_010-kuzmin-tmi_equivariant_cell_graph_transformer/runs/lzs9pcj3}{\texttt{lzs9pcj3}} & R (010 split) & \texttt{cabbi\_002} & learnable table & cosine, 30-epoch first cycle, mean & 64 & 0.452 (24) & yes, 010 eval run \\
CGT M02 & \href{https://wandb.ai/zhao-group/torchcell_010-kuzmin-tmi_equivariant_cell_graph_transformer/runs/yv4r30bi}{\texttt{yv4r30bi}} & R (010 split) & \texttt{cabbi\_002} & learnable table & cosine, 30-epoch first cycle, mean & 64 & 0.447 (25) & yes, 010 eval run \\
CGT M03 & \href{https://wandb.ai/zhao-group/torchcell_010-kuzmin-tmi_equivariant_cell_graph_transformer/runs/c7671wgj}{\texttt{c7671wgj}} & R (010 split) & \texttt{cabbi\_002} & learnable table & cosine, longer first cycle, mean & 49 & 0.462 (24) & yes, 010 eval run \\
GH 1598 & \href{https://wandb.ai/zhao-group/torchcell_025-solid-growth_equivariant_cell_graph_transformer/runs/0yw7moue}{\texttt{0yw7moue}} & R & \texttt{s0\_r\_kl\_000} & learnable table & cosine, 010 schedule, sum & 36 & 0.446 (14) & no \\
GH 1640 & \href{https://wandb.ai/zhao-group/torchcell_025-solid-growth_equivariant_cell_graph_transformer/runs/327csnlk}{\texttt{327csnlk}} & Q & \texttt{s0\_q\_kl\_004} & learnable table & cosine, 010 schedule, sum & 36 & 0.199 (7) & epoch 7 checkpoint, CPU \\
IGB 2391132 & \href{https://wandb.ai/zhao-group/torchcell_025-solid-growth_equivariant_cell_graph_transformer/runs/s1vx2zgw}{\texttt{s1vx2zgw}} & Q & \texttt{s0\_q\_kl\_emb\_017} & four-region sequence composite & constant 2.5e-4, 30 epochs, perturbed CLS & 30 & 0.262 (2) & no \\
IGB 2391133 & \href{https://wandb.ai/zhao-group/torchcell_025-solid-growth_equivariant_cell_graph_transformer/runs/8aa08xx0}{\texttt{8aa08xx0}} & Q & \texttt{s0\_q\_kl\_calm\_020} & CaLM codon embedding & constant 2.5e-4, 30 epochs, perturbed CLS & 30 & 0.252 (7) & no \\
IGB 2391134 & \href{https://wandb.ai/zhao-group/torchcell_025-solid-growth_equivariant_cell_graph_transformer/runs/pmkzwwzw}{\texttt{pmkzwwzw}} & Q & \texttt{s0\_q\_kl\_prot\_021} & ProtT5 protein embedding & constant 2.5e-4, 30 epochs, perturbed CLS & 30 & 0.270 (3) & no \\
\bottomrule
\end{tabular}

\end{table}

\notesec{What each configuration turns on}
The arm Q runs differ from one another in three settings and share everything
else. GH 1640 is the 010 configuration on the disjoint split: the learnable gene
table, the cosine schedule with a 30 epoch first cycle, sum pooling in the readout,
and the KL penalty on layer-1 attention with the corrected edge normalization. The
three IGB runs come from the joint-fitness-head branch (PR \#346, not yet on
\texttt{main}): they replace the learnable table with a fixed sequence embedding of
matched parameter count, read out through a perturbed CLS token, and train at a
constant learning rate of $2.5 \times 10^{-4}$ for 30 epochs. They share the split
and the graph penalty with GH 1640 and nothing else a replicate would have to
share, so they are four runs on one split, not one run replicated.

\notesec{The control that is missing}
The IGB runs change the gene input together with the readout and the schedule, so a
difference between any of them and GH 1640 has three candidate explanations. The
arm that separates them is the learnable table under the IGB schedule and readout,
\texttt{cgt\_s0\_q\_kl\_ctrl\_016}. Its config exists on the branch and it had no
run on W\&B as of 2026-09-15. Until it runs, the embedding variants are validation
trajectories beside GH 1640 and not a comparison with it.

%% SOURCE: experiments/025-solid-growth/scripts/additive_baselines_025_panels.py ->
%% tables/t3-disjoint-runs.tex, from the runs' W&B histories cached in
%% results/additive_baselines_025_arm_val_history.csv and
%% results/additive_baselines_025_disjoint_runs_history.csv.
\begin{table}[htbp]
\centering
\footnotesize
\caption{The four arm Q runs by their validation trajectory. Epochs is the number
logged.}
\label{tab:disjointruns}
%% GENERATED by experiments/025-solid-growth/scripts/additive_baselines_025_panels.py
%% SOURCE: the result files named in that script's docstring. Do not edit by hand.
\begin{tabular}{llp{27mm}p{33mm}rrr}
\toprule
Run & Config & Gene input & Schedule, readout & Epochs & Val max (epoch) & Val last \\
\midrule
GH 1640 & \texttt{cgt\_s0\_q\_kl\_004} & learnable table & cosine, 010 schedule, mean & 36 & 0.199 (7) & 0.131 \\
IGB 2391132 & \texttt{cgt\_s0\_q\_kl\_emb\_017} & four-region sequence composite & constant $2.5 \times 10^{-4}$, perturbed CLS & 30 & 0.262 (2) & 0.222 \\
IGB 2391133 & \texttt{cgt\_s0\_q\_kl\_calm\_020} & CaLM codon embedding & constant $2.5 \times 10^{-4}$, perturbed CLS & 30 & 0.252 (7) & 0.208 \\
IGB 2391134 & \texttt{cgt\_s0\_q\_kl\_prot\_021} & ProtT5 protein embedding & constant $2.5 \times 10^{-4}$, perturbed CLS & 30 & 0.270 (3) & 0.135 \\
\bottomrule
\end{tabular}

\end{table}

\subsection{Validation, then the one test score\secstatus{todo}}

%% SOURCE: experiments/025-solid-growth/scripts/additive_baselines_025_panels.py ->
%% additive_baselines_025_fig2_disjoint_runs.svg, from the two history caches, the
%% arm Q val and test rows of results/additive_baselines_025.csv, and
%% results/cgt_checkpoint_scores_327csnlk.json.
\tcfig{figures/additive_baselines_025_fig2_disjoint_runs.pdf}{Every transformer run
on the query-pair-disjoint split. \textbf{a}, validation Pearson per epoch on the
37{,}705 validation records of 43 held-out screens, with the additive ridge and the
three-seed MLP fit on the same validation part as horizontal lines; open circles
mark each run's best epoch. \textbf{b}, the same runs reduced to their best epoch
(open) and last logged epoch (filled). \textbf{c}, the test comparison on the
37{,}791 test records of 46 held-out screens: the one scored checkpoint, GH 1640 at
its validation-selected epoch 7, against the arm Q test nulls. The three empty
slots are the seeds of the 010 configuration that have not been run, each to be
scored at its best validation epoch.}{fig:disjointruns}

\notesec{Every run peaks early and falls}
GH 1640 passes the additive validation null of 0.150 in epoch 4, peaks at 0.199 in
epoch 7, is last above the null in epoch 16 and ends at 0.131, below every null but
B3. The three IGB runs peak higher, 0.252 to 0.270 in epochs 2 to 7, and decline
from there. The composite and CaLM embeddings end at 0.222 and 0.208 and stay above
the ridge through their last five epochs; ProtT5 ends at 0.135, level with GH 1640.
Arm R has the opposite shape, a margin over its null that appears in epoch 2 and
holds for the rest of the run (Figure~\ref{fig:ladders}c).

\notesec{The epoch 7 checkpoint on test}
The checkpoint GH 1640 saved at its best validation epoch was scored on CPU on the
arm Q test part (\file{experiments/025-solid-growth/scripts/score_cgt_checkpoint_cpu.py},
GilaHyper job 1846). The scorer reproduces the run's logged validation Pearson at
that epoch, 0.1997 recomputed against 0.1993 logged, the fp32-against-bf16 gap the
010 reproduction also showed, and two details had to be right for it to: the index
space is the genome's 6{,}607 genes, and this checkpoint pools by sum where the 010
checkpoints pooled by mean. On test the same checkpoint reaches 0.139 Pearson and
0.109 Spearman.

%% SOURCE: experiments/025-solid-growth/scripts/paired_bootstrap_025.py ->
%% tables/t6-armq-paired.tex and results/paired_bootstrap_025.json.
\begin{table}[htbp]
\centering
\footnotesize
\caption{GH 1640's epoch 7 checkpoint against three nulls on the arm Q test part,
37{,}791 records. Gain is the transformer's Pearson minus the baseline's; the
interval is a paired bootstrap over records, 2{,}000 resamples, percentile method.
Residual correlation is between $y - \hat{y}$ of the two models; prediction
correlation is between the two prediction vectors.}
\label{tab:armqpaired}
%% GENERATED by experiments/025-solid-growth/scripts/paired_bootstrap_025.py
%% SOURCE: results/paired_bootstrap_025.json. Do not edit by hand.
\begin{tabular}{lrrrrr}
\toprule
Against & Baseline $r$ & Gain & 95\% interval & Residual corr. & Prediction corr. \\
\midrule
B1 additive ridge & 0.185 & $-0.046$ & $[-0.056, -0.036]$ & 0.944 & 0.637 \\
B2 additive plus pair & 0.180 & $-0.041$ & $[-0.051, -0.030]$ & 0.943 & 0.642 \\
B5 embedding MLP, seed 0 & 0.138 & $+0.001$ & $[-0.008, +0.010]$ & 0.960 & 0.752 \\
\bottomrule
\end{tabular}

\end{table}

The checkpoint sits 0.046 below the additive ridge, with a paired interval that
excludes zero, and 0.041 below the ridge with pair offsets. Against the nonlinear
MLP the gain is 0.001 with an interval that straddles zero: on held-out screens
the transformer and the MLP on the same features are indistinguishable, and both
are below additivity. The residual correlations of 0.94 to 0.96 say the models are
wrong on the same records; the prediction correlation with B1 of 0.64 is well below
the 0.73 to 0.76 the 010 checkpoints showed on the random split, so the
transformer departs further from the additive model on unseen screens and loses
accuracy in doing so.

\subsection{The prespecified decision\secstatus{todo}}
\label{sec:decision}

A transformer result on arm Q is a claim only under three criteria, fixed before
the scoring ran. D1 asks whether the test Pearson at the validation-selected epoch
exceeds the additive ridge's. D2 asks whether the paired-bootstrap interval of the
difference excludes zero, so that a pass on D1 is more than one draw of records.
D3 asks whether D1 and D2 hold over three seeds of the configuration, so that a
pass is more than one draw of the optimizer.

%% SOURCE: experiments/025-solid-growth/scripts/paired_bootstrap_025.py ->
%% tables/t7-decision.tex and results/paired_bootstrap_025.json.
\begin{table}[htbp]
\centering
\footnotesize
\caption{The decision on the evidence available. D3 has not been run and the
verdict is provisional on it.}
\label{tab:decision}
%% GENERATED by experiments/025-solid-growth/scripts/paired_bootstrap_025.py
%% SOURCE: results/paired_bootstrap_025.json. Do not edit by hand.
\begin{tabular}{l>{\raggedright\arraybackslash}p{78mm}>{\raggedright\arraybackslash}p{40mm}l}
\toprule
Criterion & Statement & Value & Result \\
\midrule
D1 & transformer test Pearson at its validation-selected epoch exceeds B1's & 0.139 against 0.185 & \textbf{FAIL} \\
D2 & paired-bootstrap 95 percent interval of transformer minus B1 excludes zero & -0.046, [-0.056, -0.036] & \textbf{FAIL} \\
D3 & D1 and D2 hold over three seeds of the configuration & one run, one checkpoint & \textbf{NOT RUN} \\
\bottomrule
\end{tabular}

\end{table}

\textbf{Verdict, on one checkpoint: the transformer in the 010 configuration is not
shown to beat the additive null on held-out screens, and the one scored checkpoint
is below it.} D1 and D2 fail in the direction of the null, so the conclusion is not
an absence of evidence: on these 46 screens the additive ridge is the better
predictor by 0.046 Pearson with an interval that excludes zero. What the verdict
does not cover is the configuration's capability, which one run and one checkpoint
do not measure; D3 is what turns this into a statement about the model rather than
about GH 1640.

\notesec{Hypothesis, untested}
The decline after the peak in every arm Q run is consistent with the model fitting
screen-specific structure of the training pairs that does not transfer, the
quantity B4 measures. No measurement here separates that from ordinary overfitting
of label noise, and no mechanism is claimed.

\subsection{Corrections recorded\secstatus{todo}}
\label{sec:corrections}

Three claims made earlier in this analysis were wrong and are corrected here
rather than silently folded in. The Delta package for the graph-regularization
sweep carried 010's own query-pair-disjoint split, 285, 67 and 68 pairs, rather
than arm Q's, 331, 43 and 46; a result on it would have compared to neither the arm
Q nulls nor GH 1640, and the index artifact now carries arm Q's partition by query
pair name and asserts the fold sizes reproduce. The first version of this document
stated GH 1640 first cleared its validation null in epoch 3 and that GH 1598's
validation maximum lay inside the 010 checkpoints' band; the history gives epoch 4,
and 0.446 lies 0.001 below the band. And the placeholder for the test comparison
described GH 1640's checkpoint as unscorable without a GPU; the encoder runs once
and the per-record work is one cross-attention, so it scored in 13 minutes on 32
CPUs.
